% GENERATED FILE. Edit canonical lessons, then run npm run generate:books.
% canonical-source-hash: fnv1a64:091cf5afe6487b38
% canonical-lessons: ES-C33-verde, ES-C33-amarillo

\chapter{Green and Yellow}
\label{ch:spanish-green-and-yellow}
\hlchaptermodality{146}

\begin{chapteropening}
\textbf{By the end of this chapter:} \emph{I can name green and yellow in Spanish and explain why amarillo is built on the word for bitter.}

Say verde and amarillo, tie verde to virēre, to flourish, and amarillo to amarellus, a little bitter — the same root as everyday amargo.
\end{chapteropening}

\section[verde]{verde — the colour named after being alive}

\label{lesson:ES-C33-verde}



\begin{quote}
After \emph{azul}'s long journey from Afghanistan, here is a colour that stayed exactly where it started. It is also the one that names a \textbf{process} rather than a thing.
\end{quote}

\begin{cousinweb}
\textbf{verde} (\textquotedblleft{}\textbf{green}\textquotedblright{}) $\leftarrow$ Latin \textbf{viridis} (worn down through \emph{virdis}), which comes from the verb \textbf{virēre}, \textquotedblleft{}\textbf{to be green, to flourish, to be vigorous}.\textquotedblright{}

So the adjective is built on a \textbf{verb of thriving}. Latin did not name the colour of leaves; it named what leaves are \emph{doing}, and let the colour follow. The same root gives \emph{vigour} its neighbours in English: \textbf{verdant}, \textbf{verdure}, and the green-tinged \textbf{verdigris} (\textquotedblleft{}green of Greece\textquotedblright{}).

One careful non-connection: \textbf{verde is not a cousin of English \emph{green}}. \emph{Green} comes from a completely separate \textbf{Germanic} root, tied to \emph{grow} and \emph{grass} — which, funnily enough, is the same idea arrived at independently. Two families, two words, one metaphor: green is the colour of things that are growing.
\end{cousinweb}

\subsection*{Your turn}
Say these aloud:

\begin{itemize}
\raggedright
  \item \textquotedblleft{}verde\textquotedblright{} — green
  \item the root — \textquotedblleft{}virēre, to flourish\textquotedblright{} $\to$ \textquotedblleft{}verde\textquotedblright{}
  \item \textquotedblleft{}verde\textquotedblright{} then English \textquotedblleft{}verdant\textquotedblright{} — and \emph{not} \textquotedblleft{}green\textquotedblright{}
\end{itemize}

\begin{tcolorbox}[breakable,colback=teal!4,colframe=teal!35!black,title={Before you move on}]
What Latin verb gives \textbf{verde}, and what does it mean? (\textbf{Virēre}, \textquotedblleft{}to be green, to flourish.\textquotedblright{}) Is \emph{verde} related to English \emph{green}? (\textbf{No} — \emph{green} is Germanic, from the same family as \emph{grow} and \emph{grass}; the two languages reached the same metaphor separately.) Next: a colour that is not named after a colour at all.
\end{tcolorbox}
\section[amarillo]{amarillo — the colour that started as a taste}

\label{lesson:ES-C33-amarillo}



\begin{quote}
\emph{Verde} was exactly what you would expect from Latin. This one is not. Latin had three perfectly good words for yellow, and Spanish took none of them.
\end{quote}

\begin{cousinweb}
\textbf{amarillo} (\textquotedblleft{}\textbf{yellow}\textquotedblright{}) does \textbf{not} trace to \emph{flāvus}, \emph{gilvus} or \emph{lūteus} — Latin's actual yellow words. It comes from Late Latin \textbf{amarellus}, \textquotedblleft{}\textbf{a little bitter},\textquotedblright{} a diminutive of \textbf{amarus}, \textquotedblleft{}\textbf{bitter}.\textquotedblright{}

A colour named after a \textbf{taste}. The bridge is medicine: yellow was the colour of \textbf{bile} and of \textbf{jaundice} (which Spanish still calls \emph{ictericia}, but which folk speech long tied to bitterness), and of a great many bitter-tasting plants. A word meaning \textquotedblleft{}somewhat bitter\textquotedblright{} came to mean \textquotedblleft{}the yellowish look of a bitter thing,\textquotedblright{} then simply \textquotedblleft{}yellowish,\textquotedblright{} and by Spanish's Golden Age it had dropped the taste entirely and kept only the colour.

You can still check the family in a living word: \textbf{amargo} (\textquotedblleft{}bitter\textquotedblright{}) is \emph{amarillo}'s cousin, straight from the same \emph{amarus}. Say them together and the shared \emph{amar-} is right there.
\end{cousinweb}

\subsection*{Your turn}
Say these aloud:

\begin{itemize}
\raggedright
  \item \textquotedblleft{}verde, amarillo\textquotedblright{} — green and yellow
  \item \textquotedblleft{}amarillo, amargo\textquotedblright{} — hear the shared \emph{amar-}
  \item the path — \textquotedblleft{}a little bitter\textquotedblright{} $\to$ \textquotedblleft{}yellowish\textquotedblright{} $\to$ \textquotedblleft{}yellow\textquotedblright{}
\end{itemize}

\begin{tcolorbox}[breakable,colback=teal!4,colframe=teal!35!black,title={Before you move on}]
Does \textbf{amarillo} come from a Latin word for yellow? (\textbf{No} — from \emph{amarellus}, \textquotedblleft{}a little bitter,\textquotedblright{} a diminutive of \emph{amarus}.) What everyday Spanish word is its cousin? (\textbf{Amargo}, \textquotedblleft{}bitter.\textquotedblright{}) And \emph{verde}? (From \textbf{virēre}, \textquotedblleft{}to flourish\textquotedblright{} — a colour named after being alive, as this one is named after a taste.)
\end{tcolorbox}
